\section{Methodology}
\label{sec:methodology}

\para{Research question.} We ask whether induction-like heads in \gptsmall materially increase unintended copying pressure on prompts that should instead produce diffuse or random outputs. The experiment is designed to connect head-level interventions to two behaviors in the same model: intended copying on a standard induction-style continuation task and unintended copying on a constrained random-choice task.

\subsection{Model, toolkit, and compute}
\label{sec:model-setup}

We use \gptsmall with TransformerLens, which gives direct access to attention heads and head-level activation hooks. The local environment used PyTorch 2.11.0+cu130, Transformers 5.5.4, Datasets 4.8.4, SciPy 1.17.1, Pandas 3.0.2, Matplotlib 3.10.8, and Seaborn 0.13.2. The machine exposed four NVIDIA RTX A6000 GPUs through \texttt{nvidia-smi}, but PyTorch could not use CUDA because the installed driver-runtime pair was incompatible with the local torch build. All runs therefore executed on CPU. This constrained sample sizes but kept the intervention protocol identical across conditions.

\subsection{Head identification}
\label{sec:head-identification}

We identify candidate induction heads using synthetic repeated-token prompts over six single-token color candidates: \texttt{red}, \texttt{blue}, \texttt{green}, \texttt{yellow}, \texttt{black}, and \texttt{white}. For each head, we compute TransformerLens's \texttt{detect\_head(..., "induction\_head")} score across 48 head-detection prompts in the final run. The top two heads are \headpair with scores 0.740 and 0.548, respectively. \Figref{fig:head-scores} shows the resulting score profile, which places the highest-scoring heads in the familiar mid-layer region reported by prior work \cite{olsson2022induction, crosbie2024induction}.

\subsection{Tasks}
\label{sec:tasks}

\para{Copy task.} The copy task uses prompts of the form \texttt{A B C D | A B C} and evaluates whether the model predicts \texttt{D} next. This is a direct induction-style continuation pattern. We evaluate 96 prompts in the final run and measure top-1 accuracy and mean target logit.

\para{Random-choice leakage task.} The random-choice task asks the model to choose one item from the same six-item candidate set with wording such as ``with no preference'' or ``at random.'' We compare a neutral prompt against a matched prompt that additionally includes repeated distractor context, such as \texttt{Previous answers: red, red, red}. For each of 96 prompts, we measure the probability assigned to the repeated distractor in both variants.

Our primary leakage metric is the repeated-token bias delta
\begin{equation}
\Delta_{\text{bias}} = P(\text{repeated}\mid \text{biased prompt}) - P(\text{repeated}\mid \text{neutral prompt}).
\end{equation}
Higher values indicate stronger prompt-induced pressure toward the repeated distractor. We also compute normalized entropy over the six valid candidates,
\begin{equation}
H_{\text{norm}} = \frac{-\sum_{i=1}^{6} p_i \log p_i}{\log 6},
\end{equation}
where larger values indicate a more diffuse distribution, and KL divergence to the uniform distribution as a secondary diagnostic.

\para{Utility check.} To measure whether the intervention preserves general language modeling quality, we compute negative log-likelihood and perplexity on 12 non-empty samples from \wikitext using up to 64 tokens per sample.

\subsection{Interventions and baselines}
\label{sec:interventions}

We compare three conditions:
\begin{itemize}
    \item \textbf{Full model:} no intervention.
    \item \textbf{Induction-head ablation:} zero the selected heads' \texttt{hook\_z} activations in TransformerLens.
    \item \textbf{Random-head ablation control:} zero a matched number of randomly chosen heads and average across five random head sets.
\end{itemize}

The main final run uses top-2 induction-like heads, while sensitivity runs test two alternate high-scoring subsets: \texttt{L5H5,L6H9} and \texttt{L5H5,L6H9,L5H1}. The random-head control is important because hard head ablation can change model behavior for many generic reasons unrelated to induction-like circuitry.

\subsection{Data resources}
\label{sec:data}

We use the local \texttt{datasets/synthetic\_copy\_task/} resource to validate the availability of a copy-style benchmark and to ground the experimental setup. We also use the local \texttt{datasets/wikitext\_103\_v1/} resource for the utility check. The final report records 900{,}000 synthetic-copy training rows and 1{,}801{,}350 WikiText training rows in the local environment.

\subsection{Statistical analysis}
\label{sec:statistics}

The final run uses paired prompt-level comparisons for the random-choice task. Because Shapiro-Wilk tests on the paired differences rejected normality, we report Wilcoxon signed-rank tests, 95\% bootstrap confidence intervals for mean differences, and Benjamini-Hochberg adjusted $p$-values over the main comparison family. We report paired effect sizes as Cohen's $d$. This analysis focuses on two questions: whether induction-head ablation differs from the full model and whether it differs from the random-head control.
